\chapter{Conclusão}
\label{cap:conclusao}

O presente trabalho desenvolveu um infográfico interativo sobre
arquitetura de servidores, materializando, em forma de artefato
didático executável em navegador, o conteúdo das aulas 6 a 12 da
disciplina de Fundamentos de Estruturas de Computadores.

O processo passou por duas iterações. A primeira versão estabeleceu
o esqueleto visual e a interatividade básica, mas foi avaliada pelo
orientador como insuficiente em profundidade técnica. A segunda
iteração --- a que este relatório descreve --- refatorou cada
módulo para incorporar tabelas comparativas com números reais
(barramentos, latências, capacidades), fórmulas explícitas (Amdahl,
paridade XOR, síndrome de Hamming), \emph{specs} de processadores
reais (Intel Core i7-5960X, Xeon Scalable, EPYC, NVIDIA Grace
Hopper), bibliografia formal \textsc{abnt} \textsc{nbr 6023}
acessível diretamente do infográfico via citações clicáveis, e um
sexto módulo integrador (\emph{stack} completo) que sintetiza os
cinco anteriores numa única visão.

\section{Aderência aos objetivos}

Confrontados com os objetivos específicos enunciados na
Seção~1.2, observamos:

\begin{itemize}
    \item \textbf{Comparativo Desktop \textit{vs.}\ Servidor}: atendido
          com a sub-aba ``Visão Geral'' e complementado com três
          aprofundamentos (Interconexão CPU, Barramentos de E/S,
          IA \& HPC) ausentes na primeira versão.
    \item \textbf{Simulação de RAID e cálculo quantitativo}: atendido
          com o simulador original mais uma calculadora de \textsc{iops}
          e capacidade útil e a aplicação numérica da Lei de Amdahl
          ao exemplo do servidor \textsc{www}.
    \item \textbf{Visualização do código de Hamming}: atendido com a
          sub-aba ``Como funciona'', que expõe os 12 bits da palavra
          Hamming(12,8) e a tabela de cobertura dos bits de paridade.
    \item \textbf{Hierarquia com latências em escala humana}: atendido
          com tabela conversora (ciclo de CPU $\to$ 1 segundo) e
          simulador de \textit{cache hit/miss} animado.
    \item \textbf{Virtualização Tipo 1 \textit{vs.}\ Tipo 2}:
          atendido com taxonomia, anéis de proteção x86 (Ring 0 a
          Ring~$-1$) e comparação dedicada VM \textit{vs.}\ contêiner.
    \item \textbf{Integração final}: atendido com o Módulo 6
          (\emph{stack} completo), inexistente na primeira versão.
\end{itemize}

\section{Limitações}

Algumas escolhas conscientes restringem o alcance do artefato:

\begin{enumerate}
    \item Por se tratar de aplicação \emph{client-side} sem
          \emph{back-end}, nenhum dos simuladores executa medições
          reais de hardware; todos os números (\textsc{iops},
          latência, \emph{speedup}) são calculados via fórmulas
          determinísticas com parâmetros tabulares ou via amostragem
          de uma distribuição empírica.
    \item A análise da Lei de Amdahl considera apenas o cenário
          ideal de paralelização homogênea entre discos, ignorando
          contenção de barramento, \emph{cache pollution} e custo
          do controlador \textsc{raid}.
    \item A visualização do escalonador \textsc{cfs} é estática
          (árvore exemplificativa); não há simulador de inserção e
          remoção de \emph{tasks} ao longo do tempo.
    \item O suporte mobile é funcional (\emph{media queries} em três
          \emph{breakpoints}), mas tabelas extensas precisam de
          rolagem horizontal em telas estreitas.
\end{enumerate}

\section{Trabalhos futuros}

\begin{itemize}
    \item \textbf{Animar a coerência \textsc{mesi/mesif}}: implementar
          um simulador onde dois \emph{cores} disputam uma linha de
          cache e o usuário acompanha as transições de estado.
    \item \textbf{Visualização do \emph{page walk}}: animar a
          tradução de um endereço virtual nos 4 níveis da
          paginação x86-64, com e sem \textit{TLB hit}.
    \item \textbf{Integração com benchmarks reais}: incorporar
          dados de \texttt{fio}, \texttt{stress-ng} e
          \texttt{lmbench} para ancorar os valores teóricos em
          medições reproduzíveis.
    \item \textbf{Tradução para o inglês}: o infográfico tem
          potencial além do escopo da disciplina --- com tradução,
          poderia servir como recurso aberto (CC-BY) para outras
          turmas de Arquitetura de Computadores.
    \item \textbf{Acessibilidade}: auditar para conformidade
          \textsc{wcag} 2.1 \textsc{aa}, garantindo navegação por
          teclado e leitor de tela em todos os simuladores.
\end{itemize}

\section{Considerações finais}

O salto de profundidade entre as duas iterações --- de 1\,060 para
2\,554 linhas de código, de 3 para 17 referências bibliográficas,
de zero para 11 tabelas quantitativas e de 5 para 8 simuladores
interativos --- ilustra que a profundidade técnica de um artefato
didático é fundamentalmente uma decisão \emph{de conteúdo}, não de
arquitetura de software. Manter o projeto em um único arquivo HTML
sem dependências facilitou justamente esse aprofundamento: toda a
energia da segunda iteração pôde ser direcionada para a leitura
crítica do material das aulas e para a tradução dos conceitos em
visualizações executáveis, sem ser desperdiçada em configuração de
\emph{build tools} ou gerenciamento de pacotes.

O resultado é um artefato que, esperamos, faça jus à
crítica original do orientador e demonstre, de forma navegável e
auto-contida, por que um servidor é mais do que ``um computador
maior'' --- e quais subsistemas precisam estar bem compreendidos
para projetá-lo, operá-lo ou estudá-lo.
